\begin{minipage}{0.54\linewidth}
Un générateur de force électromotrice $E=\qty{24}{\V}$ et de résistance interne
$r=\qty{3}{\ohm}$ alimente un circuit électrique contenant un électrolyseur de
force contre-électromotrice $E^{\prime}$ et de résistance interne
$r^{\prime}=\qty{4}{\ohm}$.

Le générateur débite un courant électrique d'intensité~: $I=\qty{0.5}{\A}$
pendant une durée $\Delta t=\qty{30}{\min}$.
\end{minipage}
\hfill
\begin{minipage}{0.44\linewidth}
    \centering
    \begin{circuitikz}
        \draw (0,0) node[below=2pt] {$N$}
                    to[battery2,invert,l={$(E,r)$},*-*] ++(0,2.5)
                    node[above=2pt] {$P$}
                    to[short,i=$I$] ++(4,0)
                    node[above=2pt] {$A$}
                    to[dcvsource,name=E,l={$(E^{\prime},r^{\prime})$},*-*]
                    ++(0,-2.5)
                    node[below=2pt] {$B$}
                    to[short] (0,0);
        \draw[thick] (E.left) -- ++(0,-3.2mm) (E.right) -- ++(0,3.2mm);
    \end{circuitikz}
\end{minipage}
\begin{enumerate}
    \item Par application de la loi des mailles, montrer que la valeur de la
          force contre-électromotrice de l'électrolyseur est~:
          $E^{\prime}=\qty{20.5}{\V}$~\textbf{(1pt)}
    \item Calculer $W_{t}$ l'énergie totale du générateur.~\textbf{(1pt)}
    \item Calculer $W_{e}$ l'énergie électrique fournie par le
          générateur.~\textbf{(1pt)}
    \item En déduire le rendement du générateur.~\textbf{(1pt)}
    \item Calculer $W_{u}$ l'énergie utile de l'électrolyseur.~\textbf{(1pt)}
    \item Calculer $W_{J}$ l'énergie électrique dissipée par effet joule dans
          l'électrolyseur.~\textbf{(1pt)}
    \item Calculer $\rho_{c}$ le rendement global du circuit.~\textbf{(1pt)}
\end{enumerate}

